\BourbakiExercises{Exercises}

\begin{enumerate}
\def\labelenumi{\arabic{enumi})}
\tightlist
\item
  \begin{enumerate}
  \def\labelenumii{\alph{enumii})}
  \tightlist
  \item
    Let D be a noncommutative field with center K whose elements all have degree \(\leqslant 2\) over K. Prove that D is a quaternion algebra over K (Exercise 6, b) of VIII, p.~270).
  \end{enumerate}
\end{enumerate}

\begin{enumerate}
\def\labelenumi{\alph{enumi})}
\setcounter{enumi}{1}
\tightlist
\item
  Prove that a Cayley K-algebra that is a noncommutative field with center K is a quaternion algebra over K.
\end{enumerate}

¶ 2) a) Let A be a simple ring of finite rank over its center K, and let s be an involutive antiautomorphism of A. Let \(s'\) be an involutive antiautomorphism of A that coincides with s on K. Prove that there exists an element a of A such that we have \(s'(x) = as(x)a^{-1}\) for every \(x \in A\) and \(s(a) = a\) or \(s(a) = -a\) (apply the Skolem--Noether theorem and Theorem 3 of V, §11, No.~6, p.~85). Also prove the converse.

\begin{enumerate}
\def\labelenumi{\alph{enumi})}
\setcounter{enumi}{1}
\item
  Suppose that the restriction of s to K is not the identity; let \(K_{0}\) be the subfield of K consisting of the elements fixed by s. Prove that the set \(A_{0}\) (resp. \(A_{1}\) ) of elements a of A such that \(s(a) = a\) (resp. \(s(a) = -a\) ) is a \(K_{0}\) -linear subspace of A of dimension {[}A : K{]} and that \(A_{0}\) is a \(K_{0}\) -structure on A (use V, §10, No.~4, p.~63, Proposition 7).
\item
  Suppose that the restriction of s to K is the identity; set \([A:K]=m^{2}\) . Prove that \(A_{0}\) is a \(K_{0}\) -linear subspace of A of dimension \(m(m+1)/2\) or \(m(m-1)/2\) (reduce to the case when A is a matrix algebra over K and use a)).
\end{enumerate}

¶ 3) Let \((F, s)\) be a semisimple (associative) Cayley K-algebra.

\begin{enumerate}
\def\labelenumi{\alph{enumi})}
\item
  Prove that the K-algebra F is simple or isomorphic to the product of a simple algebra and its opposite algebra (consider the action of conjugation on the set of central idempotents).
\item
  Prove that if F is not commutative, then it is isomorphic to a quaternion algebra over K (use Exercise 1).
\item
  Prove that if F is commutative, then we are in one of the following cases:
\end{enumerate}

\(\alpha)\) F = K, and s is the identity.

\(\beta)\) F = K × K, and \(s(\lambda, \mu) = (\mu, \lambda)\) .

\(\gamma)\) F is a separable quadratic extension of K, and \(s\) is the nontrivial automorphism of F.

\(\delta)\) K has characteristic 2, F is an 2-radical extension of K of height 1, and \(s\) is the identity.

¶ 4) Let \((F, s)\) be an Artinian (associative) Cayley K-algebra, and let \(\mathfrak{r}\) be its radical. a) Prove that we have \(x^{2} = 0\) and \(s(x) = -x\) for every \(x \in \mathfrak{r}\) ; the algebra \(\overline{\mathrm{F}} = \mathrm{F} / \mathfrak{r}\) is Cayley and isomorphic to one of the algebras described in Exercise 3.

\begin{enumerate}
\def\labelenumi{\alph{enumi})}
\setcounter{enumi}{1}
\item
  Prove that if \(\overline{F}\) is a quaternion algebra over K, then \(\mathfrak{r}\) is reduced to 0 (observe that we have \((ab - ba)x = 0\) for \(a, b \in F, x \in \mathfrak{r}\) ).
\item
  Suppose that \(\overline{F}\) is commutative. Prove that \(\mathfrak{r}^{3}\) is reduced to 0 but that \(\mathfrak{r}^{2}\) may not be.
\item
  Suppose that \(\overline{F}\) is commutative and \(\mathfrak{r}^{2}\) is reduced to 0. Prove that there exists an \(\overline{F}\) -module V such that F is isomorphic to \(\overline{F} \oplus V\) endowed with the structure of a K-algebra for which we have \((\lambda, v)(\mu, w) = (\lambda \mu, \lambda w + s(\mu) v)\) . (In cases \(\alpha\) ) through \(\gamma\) ) of Exercise 3, c), use Corollary 1 of VIII, p.~245. In case \(\delta\) ), choose a p-basis \((e_{i})_{i \in I}\) of \(\overline{F}\) over K and for each \(i \in I\) , an element \(f_{i}\) of F that lifts \(e_{i}\) , and prove that the subalgebra of F generated by the \((f_{i})\) is isomorphic to \(\overline{F}\) .
\end{enumerate}

¶ 5) Suppose that the characteristic of K is ≠ 2. Let F be a not necessarily associative alternative Cayley K-algebra (III, Appendix, No.~2); we denote its conjugation by \(s : x \mapsto \overline{x}\) and its Cayley trace and norm by T and N, respectively. We say that a subalgebra E of F is nondegenerate if for every nonzero element x of E, there exists an \(x' \in E\) such that \(\mathrm{T}(x\overline{x}') \neq 0\) .

\begin{enumerate}
\def\labelenumi{\alph{enumi})}
\item
  Prove the equality \(x(\overline{z}y) + z(\overline{x}y) = \mathrm{T}(\overline{x}z)y = (yx)\overline{z} + (yz)\overline{x}\) for x, y, z in F (reduce to the case z = x).
\item
  Let E be a subalgebra of F containing 1, distinct from F, stable under s, and finite-dimensional over K; suppose that E and F are nondegenerate. Prove that there exists a nonzero element j of F such that \(T(xj) = 0\) for every \(x \in E\) and \(N(j) \neq 0\) . We have \(\overline{j} = -j\) , \(xj = -j\overline{x}\) for every \(x \in E\) , and \(j^{2} = \gamma1_{F}\) with \(\gamma = -N(j)\) .
\item
  Let \(x\) and \(y\) be elements of E. Prove the formulas \(x(yj) = (yx)j = (yj)\overline{x}\) and \((xj)(yj) = \gamma \overline{y}x\) (use \(a\) ) and Exercise 2 of III, Appendix, p.~654).
\end{enumerate}

\(d\) ) Prove that \(\mathrm{E} + \mathrm{E}j\) is a nondegenerate subalgebra of F, stable under \(s\) , and of dimension \(2\dim (\mathrm{E})\) that can be identified with the Cayley extension of \((\mathrm{E},s)\) defined by \(\gamma\) . Deduce from Proposition 3 of III, Appendix, No.~2, p.~614 that E is associative.

\begin{enumerate}
\def\labelenumi{\alph{enumi})}
\setcounter{enumi}{4}
\tightlist
\item
  Deduce from the above that the nondegenerate alternative Cayley K-algebras are the following:
\end{enumerate}

\begin{enumerate}
\def\labelenumi{(\roman{enumi})}
\item
  the K-algebra \(\mathrm{K}\) endowed with the identity involution
\item
  the quadratic algebras over K of type \((\alpha,0)\) with \(\alpha\neq0\)
\item
  the quaternion algebras over K of type \((\alpha,0,\gamma)\) with \(\alpha\gamma\neq0\)
\item
  the octonion algebras over \(\mathrm{K}\) of type \((\alpha, 0, \gamma, \delta)\) with \(\alpha \gamma \delta \neq 0\) .
\end{enumerate}

(Construct a sequence of Cayley extensions \(K \subset E_{1} \subset E_{2} \subset \cdots \subset F\) , and observe that an octonion algebra is not associative.)

\begin{enumerate}
\def\labelenumi{\alph{enumi})}
\setcounter{enumi}{5}
\tightlist
\item
  Prove that an alternative Cayley K-algebra whose nonzero elements are all invertible is one of the aforementioned types.
\end{enumerate}

\(g\) ) Let R a maximal ordered field. Prove that an alternative Cayley R-algebra whose nonzero elements are all invertible is isomorphic to R, C, H, or O (VIII, p.~368, Remark 1).

\begin{enumerate}
\def\labelenumi{\arabic{enumi})}
\setcounter{enumi}{5}
\tightlist
\item
  Let H be the quaternion field of type \((-1,0,-1)\) over Q.
\end{enumerate}

\begin{enumerate}
\def\labelenumi{\alph{enumi})}
\item
  An extension \(\mathrm{K}\) of \(\mathbf{Q}\) splits \(\mathrm{H}\) if and only if \(-1\) is the sum of two squares in \(\mathrm{K}\) (apply Propositions 3 of VIII, p.~363 and 4 of VIII, p.~364).
\item
  Let K be a cyclic extension of Q of degree \(2^{n}\) that splits H. Prove that no proper subfield of K splits H (observe that the unique subfield of K of degree \(2^{n-1}\) is the intersection of K and a maximal ordered field containing Q).
\item
  Let k be an integer \(\geqslant1\) and p a prime divisor of \(2^{2^{k}}+1\) . Prove that the greatest integer n such that \(2^{n}\) divides p-1 is \textgreater k (observe that we have \(2^{p-1}\equiv1\) modulo p). Let \(\Omega\) be a primitive p-the root of 1 and \(\mathrm{E}=\mathbf{Q}(\Omega)\) . Prove that -1 is the sum of two squares in E (let F be the extension \(\mathrm{E}(\sqrt{-1})\) ; show that -1 belongs to \(\mathrm{N}_{\mathrm{F}/\mathrm{E}}(\mathrm{F}^{*})\) by proving that this is true for \(\omega\) and \(1+\omega^{h}\) and by using the identity \(\prod_{j=0}^{r-1}(1+\mathrm{T}^{2^{j}})=\sum_{h=0}^{2^{r}-1}\mathrm{T}^{h})\) .
\end{enumerate}

\(d\) ) Prove that a cyclic extension of \(\mathbf{Q}\) of degree \(2^{n}\) contained in E splits H (apply VIII, p.~283, Corollary 2).

¶ 7) Suppose that the characteristic of K is different from 2.

\begin{enumerate}
\def\labelenumi{\alph{enumi})}
\item
  Let D be a field containing K and F the quaternion algebra over K of type \((\alpha, 0, \gamma)\) . Prove that \(D \otimes_{K} F\) is not a field if and only if there exist two commuting elements \(x\) and \(y\) in D such that \(x^2 - \alpha y^2 = \gamma\) (apply Exercise 4, c) of VIII, p.~351, first to \(D \otimes_{K} K(i)\) and then to \(D \otimes_{K} F\) ).
\item
  Let \(F'\) be another quaternion algebra over K. Then \(F \otimes_{K} F'\) is a field if and only if the only solution of \(\mathrm{N}_{\mathrm{F}}(q) = \mathrm{N}_{\mathrm{F}'}(q')\) , where \(q \in F\) and \(q' \in F'\) are pure quaternions (III, §2, p.~618, Exercise 3), is \(q = q' = 0\) . (Reduce to the case when F is a field and use the criterion of a). Note that if neither x nor y belongs to K, then y belongs to \(\mathrm{K}(x)\) , and write \(x = s + q\) , where \(x \in K\) and q is a pure quaternion.)
\end{enumerate}

\begin{enumerate}
\def\labelenumi{\arabic{enumi})}
\setcounter{enumi}{7}
\tightlist
\item
  Let \(K_{0}\) be a commutative field of characteristic different from 2, and let \((\mathrm{X}_{n})_{n\geqslant1}\) and \((\mathrm{Y}_{n})_{n\geqslant1}\) be two infinite sequences of variables and K be the field of rational fractions \(\mathrm{K}_{0}((\mathrm{X}_{n}),(\mathrm{Y}_{n}))\) . For any \(n\geqslant1\) , denote by \(F_{n}\) the field of quaternions over K of type \((\mathrm{X}_{n},0,\mathrm{Y}_{n})\) .
\end{enumerate}

\begin{enumerate}
\def\labelenumi{\alph{enumi})}
\item
  Prove that the tensor product \(F = \bigotimes_{n} F_{n}\) (III, §4, No.~5, p.~470) is a field with center K such that every subfield generated by a finite subset has finite rank over K (use Exercise 7, a)).
\item
  Prove that if u is an inner automorphism of F, then there exists a finite subset J of N such that the restriction of u to \(F_{n}\) , for \(n \notin J\) , is the identity. Deduce that there exists an uncountable set of noninner automorphisms of F.
\end{enumerate}

